\documentclass[../../main.tex]{subfiles}
\begin{document}

{\bf [解答]}

与えられた数列
\begin{align}
  a_1 &= 2^{N}-3 \quad (N\in\mathbb{N}_{\ge2}) \label{eq:1}\\
  a_{n+1}
  &=
  \begin{cases}
    \dfrac{a_n}{2} & (a_n\in\text{even}) \\
    \dfrac{a_n-1}{2} & (a_n\in\text{odd})
  \end{cases} \label{eq:2}
\end{align}
について考える．この$M$項の和を
\begin{align}
  A_M= \sum_{n=1}^Ma_n \label{eq:3}
\end{align}
とおき，題意の不等式
\begin{align}
  A_M \le 2^{N+1}-N-5 \label{eq:4}
\end{align}
を示す．

まず\cref{eq:2}より数列$a_n$は単調減少であり，一度$a_n=0$となればそれより大きい$n$に対して常に$a_n=0$である．
具体的に計算すると，$N=2$の時は
\begin{align}
  a_1 = 1, a_k = 0 \quad (k \ge 2=N) \label{eq:5}
\end{align}
であり，$N \ge 3$の時は
\begin{align}
  a_1 = 2^N-3, a_2 = 2^{N-1}-2, a_k=2^{N+1-k}-1 (k=3,\cdots,N), a_k = 0 (N+1 \le k) \label{eq:6}
\end{align}
となる．よっていずれの場合も，$A_M$は単調増加で $M \ge N$ の時に最大値$A_N$をとるから，
\begin{align}
  A_N \le 2^{N+1}-N-5 \label{eq:7}
\end{align}
を示せば十分である．

\subparagraph{$1^{\circ}$ $N=2$の時}

\cref{eq:5}を\cref{eq:3}に代入して
\begin{align}
  A_N = 1
\end{align}
だから\cref{eq:7}は
\begin{align}
  A_N = 1 \le 2^{2+1}-2-5 = 1
\end{align}
より成立する．

\subparagraph{$2^{\circ}$ $N \ge 3$の時}

\cref{eq:6}を\cref{eq:3}に代入して
\begin{align}
  A_N
  &=(2^N-3)+(2^{N-1}-2)+\sum_{k=3}^N(2^{N+1-k}-1)\\
  &=3\cdot2^{N-1}-5+\sum_{k=1}^{N-2}2^k-(N-2)\\
  &=3\cdot2^{N-1}-N-3+2\cdot\dfrac{2^{N-2}-1}{2-1}\\
  &=2^{N+1}-N-5
\end{align}
だから，\cref{eq:7}は成立する．
\casesend

以上から題意は示された．$\blacksquare$

\newpage

\end{document}
